\chapter{Diálogo}

\begin{objetivo}
Enfrentar os textos difíceis e ouvir como outras traduções, comentários e
tradições os leem, antes de formular a síntese.
\end{objetivo}

\section{Textos difíceis}
Textos marcados com peso \texttt{dificil} em \texttt{textos.tsv}:

\textosdotema[peso=dificil]

\section{Questões que os comentários discutem}

\paragraph{O alcance de At 15:20, 29.} Os comentários divergem sobre como
entender as quatro abstenções: como princípios de validade permanente, como
medidas para a convivência entre cristãos judeus e não judeus naquele
momento, ou com base nas leis dadas a Noé e aos estrangeiros em Israel
(\rb{Gn 9}; \rb{Lv 17–18}). Anote os argumentos de cada leitura e em que
textos se apoiam.

\paragraph{Jo 6:53–56.} Como os leitores entendem \enquote{comer a carne} e
\enquote{beber o sangue} do Filho do Homem (sentido literal ou figurado)? Que
relação têm com \rb{Lc 22:20} e com o contexto de \rb{Jo 6}?

\campovazio{Leituras diferentes}{como cada leitura entende o texto e que argumentos usa}
\campovazio{Obras consultadas}{}
